% ============================================================================
\clearpage
\section*{v4.9.0-r41 Source-Only: Distributing the Roles of the Quaternion Factor; Hypercharge and Electric Charge}
\addcontentsline{toc}{section}{v4.9.0-r41 Source-Only: Distributing the Roles of the Quaternion Factor; Hypercharge and Electric Charge}
\label{sec:v490-r41-quaternion-roles}
% ============================================================================

\paragraph{Purpose.}
The inventory of the standard-model features found the quaternion factor overloaded: it carries the Lorentz spinor
(Theorem~\ref{thm:three-generations}), the weak \(SU(2)\) (Theorem~\ref{thm:SM}) and, in v4.9.0-r33, the generation index.  Layers
v4.9.0-r37 to r40 found a carrier for the generations in triality, a separation of frames and generations by zoom level, and two
\(\mathfrak{su}(2)\) commuting with colour, from the left and from the right of \(\HH\).  This layer records the pre-registered test H26
(note \texttt{llm/Vorab\_Notiz\_H26\_Verteilung\_H\_Faktor.md}, commit \texttt{1876b55}, fixed before computing) of a distribution D of
these roles.  The charge pattern of (b) and (c) below was seen analytically before the note and is disclosed there; the computation
checks it and is not a prediction.  Earlier layers are not rewritten.

\noindent\textbf{Source-only status.}
This v4.9.0-r41 layer is source-only.  It does not build a PDF, does not assign a new DOI,
and does not authorize a public release.  The v4.9.0 Version DOI
\href{https://doi.org/10.5281/zenodo.23534954}{10.5281/zenodo.23534954} and the Concept DOI
\href{https://doi.org/10.5281/zenodo.19210728}{10.5281/zenodo.19210728} are unchanged by it.

\subsection*{v4.9.0-r41 distribution D}
\label{subsec:v490-r41-distribution}

Lorentz \(SL(2,\mathbb C)\) by left multiplication \(L_q\) (rotations) and \(iL_q\) (boosts) on \(\mathbb C\otimes\HH\); weak \(SU(2)\) by
right multiplication \(R_q\) on \(\HH\), with \(T_3=-\tfrac{i}{2}R_k\); colour \(SU(3)_c\) on \(\OO\); frames \(D_4\) on \(\mathbb C\); generations on
the triality slots of \(\OO^3\).  Hypercharge \(Y/2:=Q-\tfrac12\) with \(Q\) of v4.9.0-r34/r35; electric charge \(Q_{\rm em}:=T_3+Y/2\).  Here \(i\)
is the outer unit.

\subsection*{v4.9.0-r41 results}
\label{subsec:v490-r41-results}

\paragraph{Theoremlet v4.9.0-r41.1.}
\begin{enumerate}[leftmargin=2em]
\item[(a)] On \(\mathbb C\otimes\HH\otimes\OO\) the weak \(SU(2)\) commutes with the Lorentz group, with colour and with \(Q\).  The Lorentz
rotations are fixed by \(T\); the boosts are reversed by \(T\).
\item[(b)] The joint eigenspaces of \(Q\) and \(T_3\) carry \(Q_{\rm em}\) as follows:
\begin{center}
\small
\begin{tabular}{@{}llll@{}}
\toprule
\(Q\)-sector (colour) & \(Y/2\) & \(Q_{\rm em}\) at \(T_3=+\tfrac12,-\tfrac12\) & reading\\
\midrule
\(0\) (singlet) & \(-\tfrac12\) & \(0,\ -1\) & lepton doublet \((\nu,e)\)\\
\(\tfrac23\) (triplet) & \(+\tfrac16\) & \(\tfrac23,\ -\tfrac13\) & quark doublet \((u,d)\)\\
\(\tfrac13\) (conjugate triplet) & \(-\tfrac16\) & \(\tfrac13,\ -\tfrac23\) & antiquark doublet \((\bar d,\bar u)\)\\
\(1\) (singlet) & \(+\tfrac12\) & \(1,\ 0\) & antilepton doublet \((e^+,\bar\nu)\)\\
\bottomrule
\end{tabular}
\end{center}
\item[(c)] \(TQ_{\rm em}T^{-1}=-Q_{\rm em}\) and \(T(Y/2)T^{-1}=-Y/2\).
\item[(d)] There is no weak singlet: \(T_1^2+T_2^2+T_3^2=\tfrac34\) on the whole space.
\item[(e)] On \(\mathbb C\otimes\HH\otimes\OO^3\) the cyclic shift of the slots commutes with \(L_q\), \(R_q\), colour and \(Q\).
\end{enumerate}

\paragraph{Proof sketch.}
By explicit complex matrices on \(\HH\otimes\OO\cong\mathbb C^{32}\), with the \(\mathbb C\) factor as scalar and \(T\) as complex conjugation
(script \texttt{llm/k82\_test\_h26.py}, deviations at most \(10^{-14}\)).  (c) follows from \(TQT^{-1}=1-Q\) (v4.9.0-r34.1(d)) and
\(T(-\tfrac i2R_k)T^{-1}=\tfrac i2R_k\).  (d): \(R_q\) acts on \(\HH\cong\mathbb C^2\otimes\mathbb C^2\) as a doublet on the right index.  \(\square\)

\paragraph{Verdict and reading.}
H26a--e are confirmed; (d) is negative for D.
\begin{enumerate}[leftmargin=2em]
\item In D the quaternion factor is no longer overloaded: the Lorentz group sits on the left (front side), the weak \(SU(2)\) on the
right (back side), the generations on the triality slots.  The zoom levels separate frames from generations; the sides separate the
Lorentz group from the weak \(SU(2)\).
\item Under the choice \(Y/2=Q-\tfrac12\), \(Q_{\rm em}=T_3+Y/2\) gives exactly the hypercharges and charges of the left-handed
standard-model doublets and their conjugates, with the right colour: leptons singlets, quarks triplets, antiquarks conjugate
triplets.  The choice is motivated by the pattern and not forced by the structure.  The sector assignment is Furey's; similar patterns
are known in the literature on \(\mathbb C\otimes\HH\otimes\OO\).  New here is the placement on the sides and frames of the Companion.
\item In D the frame change \(T\) is charge conjugation, \(Q_{\rm em}\mapsto-Q_{\rm em}\).  The shift in \(Q\mapsto1-Q\) of v4.9.0-r34.1(d) is
the hypercharge offset \(\tfrac12\).  Matter frame and tachyon frame behave in D as particle and antiparticle.
\item D is incomplete.  Every state is a weak doublet; the right-handed singlets \(u_R,d_R,e_R,\nu_R\), half of a generation, are absent,
and all states carry the same left Lorentz index.  Chirality is not explained.  Weak singlets would need a structure on which the right
action of \(\HH\) is trivial on part of the space, for instance a projector coupling left and right.
\end{enumerate}

\subsection*{v4.9.0-r41 status in the three frames}
\label{subsec:v490-r41-frames}

\begin{center}
\small
\begin{tabular}{@{}>{\raggedright\arraybackslash}p{0.30\textwidth}>{\raggedright\arraybackslash}p{0.62\textwidth}@{}}
\toprule
result & M / T / P\\
\midrule
r41.1(a), (d), (e) & the same in M, T and P, except that \(T\) reverses the boosts\\
r41.1(b) & stated in M; in T all charges and hypercharges change sign (r41.1(c)): the T frame reads the conjugate states\\
r41.1(c) & relates M and T: \(T\) acts as charge conjugation; P not evaluated\\
\bottomrule
\end{tabular}
\end{center}

\subsection*{v4.9.0-r41 register entries}
\label{subsec:v490-r41-register}

\begin{center}
\small
\begin{tabular}{@{}>{\raggedright\arraybackslash}p{0.34\textwidth}>{\raggedright\arraybackslash}p{0.58\textwidth}@{}}
\toprule
statement & status after v4.9.0-r41\\
\midrule
overload of \(\HH\) & resolved in D by sides (Lorentz left, weak right) and by triality for the generations (r41.1(a), (e))\\
hypercharge and \(Q=T_3+Y/2\) & doublet pattern reproduced under the choice \(Y/2=Q-\tfrac12\) (r41.1(b)); not derived\\
frame \(T\) & charge conjugation in D (r41.1(c))\\
weak singlets, chirality & absent in D (r41.1(d)); open\\
\bottomrule
\end{tabular}
\end{center}

\subsection*{v4.9.0-r41 ledger and guard update}
\label{subsec:v490-r41-ledger-guard}

\[
R_{\rm qh}^{v490\text{-}r41}=(1,1,0)
\]
for (overload of \(\HH\) resolved in D; doublet charges reproduced under a choice; singlets and chirality absent).

The layer spends no guard:
\[
R_{\rm glob}^{(v490\text{-}r41,+)}=(0,1,1,1,1),
\qquad
\Gamma_{\rm guard}^{v490\text{-}r41}=(1,1,1,1,1),
\qquad
\beta_4=1/256\ \text{unspent}.
\]
It does not imply exact finite \(C5\) zero closure, exact all-readout terminality,
\(C6/T2\), a global photon-clock/RH theorem, or global all-readout closure.
